\documentclass[12pt,a4paper,oneside]{report}

\usepackage{iftex}
\ifPDFTeX
  \usepackage[utf8]{inputenc}
  \usepackage[T1]{fontenc}
  \usepackage{lmodern}
\else
  \usepackage{fontspec}
\fi
\usepackage[slovak]{babel}
\usepackage[a4paper,left=3cm,right=2.5cm,top=2.5cm,bottom=2.5cm]{geometry}
\usepackage{setspace}
\onehalfspacing

\usepackage{amsmath,amssymb}
\usepackage{graphicx}
\usepackage{booktabs,tabularx,multirow}
\usepackage{caption}
\usepackage{subcaption}
\usepackage{float}
\usepackage{enumitem}
\usepackage{xcolor}
\usepackage{listings}
\usepackage{tikz}
\usetikzlibrary{arrows.meta,positioning,fit,shapes.geometric,calc,babel}
\usepackage{csquotes}
\usepackage[hidelinks]{hyperref}

\graphicspath{{figures/}}
\setlength{\emergencystretch}{3em}

\newcommand{\todo}[1]{\textcolor{red}{\textbf{[TODO:} #1\textbf{]}}}
\newcommand{\optfigure}[3]{%
  \IfFileExists{figures/#1}{\includegraphics[width=#2]{#1}}{%
    \fbox{\parbox{#2}{\centering\vspace{2cm}\textcolor{red}{Chýba \texttt{figures/#1}}\\[2pt]\small #3\vspace{2cm}}}}}

\lstset{
  basicstyle=\ttfamily\small,
  breaklines=true,
  frame=single,
  columns=fullflexible,
  keepspaces=true,
  showstringspaces=false,
  literate={á}{{\'a}}1 {é}{{\'e}}1 {í}{{\'i}}1 {ó}{{\'o}}1 {ú}{{\'u}}1 {ý}{{\'y}}1
           {č}{{\v{c}}}1 {ď}{{\v{d}}}1 {ľ}{{\v{l}}}1 {ň}{{\v{n}}}1 {š}{{\v{s}}}1
           {ť}{{\v{t}}}1 {ž}{{\v{z}}}1 {ô}{{\^o}}1 {ä}{{\"a}}1 {→}{{$\rightarrow$}}1
}

\addto\captionsslovak{\renewcommand{\bibname}{Zoznam použitej literatúry}}

\begin{document}

\begin{titlepage}
  \centering
  {\large TECHNICKÁ UNIVERZITA V KOŠICIACH\par}
  {\large FAKULTA ELEKTROTECHNIKY A INFORMATIKY\par}
  \vspace{4cm}
  {\LARGE\bfseries Biometrická autentifikácia tváre\\[4pt] s ochranou proti falšovaniu\par}
  \vspace{1cm}
  {\large Semestrálne zadanie z predmetu Biometrické systémy bezpečnosti\par}
  \vfill
  \begin{flushleft}
    \begin{tabular}{@{}ll}
      Vypracovali: & Oleh Poiasnik, \todo{meno spolužiaka}\\
      Vyučujúci:   & doc.~Ing.~Matúš Pleva, PhD., Ing.~Eva Kupcová, PhD.\\
      Akademický rok: & 2026/2027\\
    \end{tabular}
  \end{flushleft}
  \vspace{1cm}
  {Košice \the\year\par}
\end{titlepage}

\chapter*{Abstrakt}
\addcontentsline{toc}{chapter}{Abstrakt}

Práca sa zaoberá prihlasovaním pomocou tváre a ochranou takého prihlasovania pred
podvrhmi, ako sú fotografie alebo videá prehrávané na displeji. Navrhli a implementovali
sme webový systém, ktorý tvár sníma bežnou webkamerou a o~výsledku rozhoduje na serveri.
Ochrana je viacvrstvová: kontrola kvality snímky, pasívna detekcia podvrhu dvojicou
neurónových sietí MiniFASNet, náhodná výzva na pohyb hlavy a kontrola, že sa osoba počas
snímania nezmenila. Tvár sa porovnáva modelom SFace. Na databáze LFW dosiahlo
rozpoznávanie EER 2,47~\% a pri použitom prahu nulové FAR pri FRR 3,90~\%. Pasívna
detekcia mala na snímkach z~databázy CASIA-FASD EER 5,25~\% a pri zvolenom prahu
APCER najviac 1,68~\%. Všetkých 30 pokusov o~prihlásenie statickým obrázkom tváre bolo
zamietnutých. Spracovanie jednej relácie trvá na bežnom procesore približne 0,3~s.

\vspace{1em}
\noindent\textbf{Kľúčové slová:} biometria, rozpoznávanie tváre, detekcia prezentačných
útokov, anti-spoofing, živosť, ISO/IEC~30107-3

\tableofcontents
\listoffigures
\listoftables

\chapter{Úvod}

Odomknúť telefón pohľadom do kamery je dnes úplne bežné a podobne sa tvárou
prihlasujeme aj do bankových aplikácií či pri overovaní totožnosti na diaľku. Je to
pohodlné -- netreba si pamätať heslo a tvár so sebou nosíme stále. Má to však háčik:
tvár nie je tajomstvo. Fotku človeka nájdeme na sociálnych sieťach za pár sekúnd a ak
ju systém nedokáže odlíšiť od skutočnej osoby, útočník sa prihlási namiesto nej.
Správa agentúry ENISA \cite{enisa2022remote} uvádza takéto podvrhy -- fotografie,
videá prehrávané na displeji či masky -- medzi hlavnými hrozbami vzdialeného overovania
totožnosti.

Samotné porovnanie dvoch tvárí teda nestačí. Systém musí okrem otázky \enquote{je to
tá istá osoba?} zodpovedať aj otázku \enquote{je pred kamerou naozaj živý človek?}.
Práve tejto oblasti, detekcii prezentačných útokov (angl.\ \emph{presentation attack
detection}, \emph{anti-spoofing}), sme sa v~projekte venovali.

Cieľom práce bolo navrhnúť, naprogramovať a otestovať webový systém, ktorý
používateľa prihlási iba pomocou tváre z~bežnej webkamery a zároveň sa nedá
oklamať fotografiou, displejom ani vopred nahraným videom. Nechceli sme sa spoľahnúť
na jedinú techniku, preto systém kombinuje viac vrstiev ochrany: kontrolu kvality
snímky, neurónovú sieť na pasívnu detekciu podvrhu, náhodnú výzvu na pohyb hlavy
a kontrolu, že sa osoba počas snímania nezmenila. Hotový systém je verejne nasadený
a jeho zdrojový kód je dostupný na GitHube.

Práca je rozdelená nasledovne. Kapitola~\ref{ch:teoria} stručne vysvetľuje, ako
funguje rozpoznávanie tváre, aké útoky naň existujú a ako sa meria úspešnosť
biometrických systémov. V~kapitole~\ref{ch:navrh} opisujeme návrh riešenia --
architektúru, spracovanie snímok a bezpečnostné opatrenia. Kapitola~\ref{ch:implementacia}
ukazuje, ako sme systém implementovali. V~kapitole~\ref{ch:testovanie} sú výsledky
testovania a vyhodnotenia, kapitola~\ref{ch:prirucka} obsahuje príručku k~inštalácii
a používaniu a na záver zhrnieme dosiahnuté výsledky a možné rozšírenia.

\chapter{Analýza problematiky}
\label{ch:teoria}

\section{Ako funguje rozpoznávanie tváre}

Rozpoznávanie tváre prebieha v~niekoľkých krokoch. Najprv treba tvár v~obraze nájsť.
Detektor vráti jej polohu a niekoľko orientačných bodov -- typicky stredy očí, špičku
nosa a kútiky úst. Moderné detektory sú pritom veľmi malé; YuNet má menej ako sto tisíc
parametrov a na bežnom procesore zvládne detekciu za niekoľko milisekúnd
\cite{wu2023yunet}. Podľa orientačných bodov sa tvár zarovná do jednotnej polohy
a veľkosti, aby porovnanie neovplyvnilo, či je hlava mierne naklonená alebo ďalej od
kamery.

Zarovnaný výrez tváre spracuje konvolučná neurónová sieť a vytvorí z~neho vektor
príznakov, akýsi číselný \enquote{odtlačok} tváre. Dobrý vektor má dve vlastnosti:
dve fotky toho istého človeka dávajú podobné vektory aj pri inom svetle či výraze
a fotky rôznych ľudí dávajú vektory, ktoré sú od seba ďaleko. Podobnosť sa najčastejšie
meria kosínusom uhla medzi vektormi a o~zhode rozhodne prah. Systém pritom môže pracovať
v~dvoch režimoch. Pri \emph{verifikácii} (1:1) používateľ tvrdí, kto je, a systém to iba
potvrdí porovnaním s~jedným uloženým vzorom. Pri \emph{identifikácii} (1:N) systém
hľadá najpodobnejší vzor v~celej databáze, takže s~rastúcim počtom používateľov rastie
aj riziko zámeny.

To, ako dobre vektory oddeľujú ľudí, určuje hlavne spôsob trénovania siete. Metóda
ArcFace pri trénovaní pridáva medzi triedy (osoby) uhlovú rezervu, takže sieť sa musí
naučiť vektory jednej osoby zhlukovať tesnejšie, než by bolo inak potrebné
\cite{deng2022arcface}. AdaFace túto rezervu mení podľa kvality obrázka -- pri
rozmazaných alebo tmavých snímkach je miernejšia, aby sa sieť neučila z~nekvalitných
dát zlé vzory \cite{kim2022adaface}. SFace, ktorý používame v~našom systéme, rieši
podobný problém tak, že silu, ktorou sa vektory priťahujú a odpudzujú, obmedzuje
sigmoidnou funkciou \cite{zhong2021sface}.

\section{Kvalita biometrickej vzorky}

Aj najlepší model zlyhá, ak dostane zlú snímku. Rozmazaná, tmavá, príliš malá alebo
silno otočená tvár zvyšuje počet odmietnutých oprávnených používateľov a niekedy sa
tvár nepodarí ani nájsť. Kvalitu sa preto oplatí kontrolovať ešte pred porovnaním
a používateľovi povedať, čo má zlepšiť. Prehľad metód hodnotenia kvality
\cite{schlett2022fiqa} chápe kvalitu ako užitočnosť snímky pre rozpoznávanie, nie ako
jej estetiku -- pekná, ale z~profilu odfotená tvár je pre systém horšia ako obyčajná
frontálna fotka z~webkamery. Norma ISO/IEC 29794-5 z~roku 2025 stanovuje, ako merať
jednotlivé vlastnosti fotografie tváre, napríklad ostrosť, osvetlenie, polohu hlavy či
otvorené oči, a rozlišuje použitie pri fotografiách do dokladov, pri registrácii a pri
bežnom overovaní \cite{iso29794-5}.

\section{Prezentačné útoky}

Pri prezentačnom útoku útočník ukáže snímaču podvrh, aby ho systém považoval za inú
osobu, alebo aby skryl vlastnú identitu. Prostriedok útoku sa označuje ako PAI
(\emph{presentation attack instrument}). Najjednoduchší je vytlačená fotografia,
prípadne ohnutá alebo s~vystrihnutými očami, aby pôsobila živšie. Častý je aj
\emph{replay} útok, pri ktorom útočník ukáže fotografiu alebo video na displeji mobilu
či tabletu. Náročnejšie sú papierové, silikónové alebo 3D tlačené masky
\cite{hernandezortega2023pad}. Pri hodnotení systému sa každý druh útoku sleduje
zvlášť, pretože proti každému pomáhajú iné vlastnosti obrazu.

Od prezentačných útokov treba odlíšiť injekčné útoky. Pri nich útočník kameru úplne
obíde a pošle do systému podvrhnutý obraz priamo, napríklad cez virtuálnu kameru alebo
upravenú aplikáciu. Takýto obraz nemusí mať žiadne fyzické stopy podvrhu, preto ho
detekcia prezentačných útokov sama nezachytí. Agentúra ENISA ich preto uvádza ako
samostatnú kategóriu hrozieb a odporúča ich riešiť ďalšími opatreniami, napríklad
overením integrity zariadenia \cite{enisa2022remote}.

\section{Detekcia prezentačných útokov}

Metódy detekcie sa delia na pasívne a aktívne. Pasívne metódy analyzujú samotný obraz
a hľadajú stopy podvrhu. Starší prístup využíval ručne navrhnuté príznaky textúry,
dnes sa takmer výhradne používajú hlboké neurónové siete, ktoré sa tieto stopy naučia
z~dát -- napríklad textúru papiera, moaré vzniknuté snímaním displeja, odlesky alebo
okraje fotografie \cite{yu2023fassurvey}. Ich výhodou je, že používateľ nemusí nič
robiť. Nevýhodou je, že model vie dobre rozpoznať iba to, čo videl pri trénovaní.

Práve generalizácia je hlavným otvoreným problémom. Model natrénovaný na jednom datasete
často zlyhá pri inej kamere, inom osvetlení alebo novom type útoku. Výskum sa preto
zameriava na metódy, ktoré lepšie oddelia živé tváre od podvrhov aj v~nových podmienkach
\cite{sun2023rethinking}, alebo využívajú veľké predtrénované modely spájajúce obraz
a text \cite{srivatsan2023flip}. Súťaže na rozsiahlych dátach z~reálneho sveta zároveň
ukazujú, že výsledky z~laboratórnych datasetov sa do praxe prenášajú len čiastočne
\cite{wang2023wildfas}.

Aktívne metódy naopak žiadajú od používateľa spoluprácu -- má žmurknúť, usmiať sa alebo
otočiť hlavu určeným smerom (tzv.\ \emph{challenge--response}). Statická fotografia
takúto akciu nevykoná a ak je poradie akcií náhodné, nepomôže ani vopred nahrané video.
Cenou je o~niečo dlhšie a menej pohodlné prihlásenie. V~praxi sa preto oba prístupy
často kombinujú, čo je aj prístup nášho systému.

\section{Metriky a normy hodnotenia}

Presnosť rozpoznávania opisujú dve chyby: FAR (podiel prijatých cudzích osôb) a FRR
(podiel odmietnutých oprávnených používateľov). Obe závisia od prahu -- prísnejší prah
znižuje FAR, ale zvyšuje FRR. Systémy sa preto porovnávajú pomocou EER, teda hodnoty,
pri ktorej sú obe chyby rovnaké, alebo pomocou krivky DET, ktorá ukazuje celý kompromis
naraz. Pre bezpečnostné aplikácie sa často uvádza aj FRR pri pevnom, veľmi nízkom FAR.

Pre detekciu prezentačných útokov definuje metriky norma ISO/IEC 30107-3
\cite{iso30107-3}. APCER je podiel útokov, ktoré systém nesprávne prijal ako živú tvár,
a počíta sa zvlášť pre každý druh útoku; rozhodujúci je najhorší z~nich. BPCER je podiel
skutočných používateľov, ktorých systém nesprávne označil za útok. V~literatúre sa
často uvádza aj ACER, priemer APCER a BPCER, no norma takéto priemerovanie neodporúča,
pretože skrýva, či systém zlyháva v~bezpečnosti alebo v~pohodlí \cite{busch2023standards}.
Vhodnejšie je uvádzať BPCER pri pevnej hodnote APCER, napríklad BPCER20 pri APCER 5~\%.

\section{Ochrana biometrických údajov}

Biometrické údaje, ktoré slúžia na jednoznačnú identifikáciu osoby, patria podľa GDPR
medzi osobitné kategórie osobných údajov a ich spracovanie vyžaduje výslovný súhlas.
Na rozdiel od hesla sa tvár po úniku nedá zmeniť, preto treba chrániť aj uložené vzory.
Od ochrany sa očakáva, že z~uloženého vzoru nepôjde zrekonštruovať pôvodnú tvár, že
vzory z~rôznych systémov sa nebudú dať prepojiť a že kompromitovaný vzor pôjde nahradiť
novým. Používajú sa na to napríklad zrušiteľná biometria, biometrické kryptosystémy
alebo homomorfné šifrovanie \cite{abdullahi2024template}. Od roku 2024 platí v~EÚ aj
akt o~umelej inteligencii, ktorý systémy diaľkovej biometrickej identifikácie zaraďuje
medzi vysokorizikové. Výnimkou je biometrická verifikácia, ktorej jediným účelom je
potvrdiť, že osoba je tá, za ktorú sa vydáva \cite{euaiact2024}.

\chapter{Návrh riešenia}
\label{ch:navrh}

Cieľom praktickej časti bolo vytvoriť funkčný webový systém, ktorý používateľa
autentifikuje iba pomocou tváre snímanej bežnou webkamerou a ktorý zároveň odolá
najčastejším prezentačným útokom -- fotografii, displeju s~fotografiou alebo videom
a pokusu o~zámenu osoby počas snímania. Kapitola opisuje architektúru, biometrický
reťazec a bezpečnostný návrh. Implementačné detaily nasledujú
v~kapitole~\ref{ch:implementacia}.

\section{Architektúra}

Systém má klient-server architektúru (obr.~\ref{fig:architektura}). Klient je
jednostránková aplikácia (SPA) v~jazyku TypeScript s~knižnicou React, ktorá ovláda
webkameru cez rozhranie \texttt{getUserMedia}, zobrazuje pokyny a odosiela snímky na
server. Server je napísaný v~jazyku Python s~rámcom FastAPI a obsahuje celý
biometrický reťazec. Všetky rozhodnutia o~živosti a zhode robí server -- kód
v~prehliadači môže útočník ľubovoľne upraviť, preto mu nemožno dôverovať. Trvalé dáta
sú v~relačnej databáze (SQLite pri vývoji, PostgreSQL v~produkcii), krátkodobý stav
(jednorazové výzvy, počítadlá pokusov) v~pamäti procesu alebo v~Redise.

\begin{figure}[H]
  \centering
  \begin{tikzpicture}[
      node distance=6mm and 8mm,
      box/.style={draw, rounded corners, align=center, minimum height=9mm,
                  minimum width=26mm, font=\small, fill=white},
      grp/.style={draw, dashed, rounded corners, inner sep=3mm},
      arr/.style={-{Stealth[length=2mm]}, thick}]
    \node[box] (cam) {Webkamera};
    \node[box, below=of cam] (spa) {React SPA\\riadenie výzvy};
    \node[grp, fit=(cam)(spa), label={[font=\small\bfseries]above:Prehliadač}] (client) {};
    \node[box, right=22mm of spa] (api) {REST API\\(FastAPI)};
    \node[box, right=of api] (svc) {Služby\\(registrácia,\\overenie)};
    \node[box, above=of svc] (eng) {Biometrický\\engine};
    \node[box, right=of eng] (models) {YuNet, SFace,\\MiniFASNet\\(ONNX)};
    \node[box, below=of svc] (db) {Databáza\\(vzory, audit)};
    \node[box, right=of svc] (store) {Výzvy,\\rate-limit};
    \node[grp, fit=(api)(svc)(eng)(models)(db)(store),
          label={[font=\small\bfseries]above:Server}] (server) {};
    \draw[arr] (cam) -- (spa);
    \draw[arr] (spa) -- node[above, font=\scriptsize]{HTTPS/JSON} (api);
    \draw[arr] (api) -- (svc);
    \draw[arr] (svc) -- (eng);
    \draw[arr] (eng) -- (models);
    \draw[arr] (svc) -- (db);
    \draw[arr] (svc) -- (store);
  \end{tikzpicture}
  \caption{Architektúra systému}
  \label{fig:architektura}
\end{figure}

Serverová časť je rozdelená do vrstiev (tab.~\ref{tab:vrstvy}). Biometrické komponenty
sú definované cez rozhrania -- detektor, extraktor príznakov a pasívny detektor živosti
možno vymeniť bez zmeny zvyšku systému a v~testoch ich nahradiť deterministickými
náhradami. Všetky objekty sa skladajú na jednom mieste aplikácie.

\begin{table}[H]
  \centering
  \caption{Vrstvy serverovej časti}
  \label{tab:vrstvy}
  \begin{tabularx}{\textwidth}{@{}l l X@{}}
    \toprule
    Vrstva & Balík & Zodpovednosť \\
    \midrule
    Prezentačná & \texttt{app/api} & HTTP rozhranie, validácia vstupov, autentifikácia JWT, mapovanie chýb \\
    Aplikačná & \texttt{app/services} & prípady použitia, transakcie, audit, dočasné zablokovanie účtu \\
    Doménová & \texttt{app/biometrics} & detekcia, kvalita, živosť, príznaky, rozhodovanie \\
    Infraštruktúrna & \texttt{app/db}, \texttt{app/core} & databáza, šifrovanie, konfigurácia, Redis \\
    \bottomrule
  \end{tabularx}
\end{table}

\section{Biometrický reťazec}

Každá snímka prechádza rovnakým spracovaním (obr.~\ref{fig:retazec}). Najprv detektor
YuNet \cite{wu2023yunet} nájde tváre a päť orientačných bodov (stredy očí, špička nosa,
kútiky úst). Z~orientačných bodov sa hneď vypočíta poloha hlavy -- aj vtedy, keď snímka
nesplní požiadavky na kvalitu, pretože počas otáčania hlavy je obraz často mierne
rozmazaný. Ak snímka kvalitu spĺňa, model SFace \cite{zhong2021sface} z~nej vytvorí
128-rozmerný vektor príznakov a pri frontálnych snímkach dvojica modelov MiniFASNet
odhadne pravdepodobnosť, že ide o~živú tvár.

\begin{figure}[H]
  \centering
  \begin{tikzpicture}[
      node distance=5mm,
      stage/.style={draw, rounded corners, align=center, font=\small,
                   minimum width=24mm, minimum height=8mm},
      arr/.style={-{Stealth[length=2mm]}, thick}]
    \node[stage] (in) {Snímka};
    \node[stage, right=of in] (det) {Detekcia\\YuNet};
    \node[stage, above right=2mm and 8mm of det] (pose) {Poloha hlavy};
    \node[stage, below right=2mm and 8mm of det] (q) {Kvalita\\vzorky};
    \node[stage, above right=1mm and 8mm of q] (emb) {SFace\\(príznaky)};
    \node[stage, below right=1mm and 8mm of q] (pad) {MiniFASNet\\(len frontálne)};
    \draw[arr] (in) -- (det);
    \draw[arr] (det) -- (pose.west);
    \draw[arr] (det) -- (q.west);
    \draw[arr] (q) -- (emb.west);
    \draw[arr] (q) -- (pad.west);
  \end{tikzpicture}
  \caption{Spracovanie jednej snímky}
  \label{fig:retazec}
\end{figure}

Rozhodnutie sa však nerobí na základe jednej snímky, ale celej \emph{relácie}.
Relácia je postupnosť snímok zachytených počas jednej náhodnej výzvy. Server ju
vyhodnotí v~piatich krokoch a pri prvom neúspešnom kroku ju zamietne s~konkrétnym
dôvodom:

\begin{enumerate}
  \item \textbf{Kvalita} -- relácia musí obsahovať aspoň jednu použiteľnú frontálnu snímku
        (dôvod \texttt{quality}).
  \item \textbf{Pasívna detekcia} -- priemerné skóre živosti frontálnych snímok musí byť
        aspoň $\tau_{\mathrm{PAD}} = 0{,}50$ a žiadna z~nich nesmie klesnúť pod $0{,}20$
        (dôvod \texttt{passive\_spoof}).
  \item \textbf{Aktívna výzva} -- v~každom kroku výzvy musí byť zaznamenaný požadovaný
        pohyb hlavy (dôvod \texttt{active\_challenge\_failed}).
  \item \textbf{Konzistencia identity} -- každá snímka relácie musí byť podobná vzoru
        relácie aspoň na $0{,}30$ (dôvod \texttt{identity\_inconsistent}).
  \item \textbf{Porovnanie} -- vzor relácie sa porovná s~uloženým vzorom (1:1) alebo
        s~galériou (1:N); zhoda nastane pri kosínusovej podobnosti
        $s \geq \tau_{\mathrm{match}} = 0{,}40$ (dôvod \texttt{no\_match}).
\end{enumerate}

Pasívna detekcia sa počíta iba na frontálnych snímkach, pretože modely MiniFASNet boli
trénované na tvárach otočených do kamery a pri otočenej hlave dávajú zbytočne nízke
skóre. Prah $0{,}50$ sme zvolili na základe merania v~kapitole~\ref{ch:testovanie}.
Vzor relácie je normalizovaný priemer vektorov príznakov frontálnych snímok:
\begin{equation}
  \mathbf{t} = \frac{\sum_{i} \hat{\mathbf{e}}_i}{\left\lVert \sum_{i} \hat{\mathbf{e}}_i \right\rVert},
  \qquad
  s(\mathbf{t}, \mathbf{r}) = \mathbf{t}^{\top}\hat{\mathbf{r}},
\end{equation}
kde $\hat{\mathbf{e}}_i$ sú L2-normalizované vektory príznakov a $\hat{\mathbf{r}}$ je
normalizovaný referenčný vzor. Prah $0{,}40$ je mierne prísnejší ako hodnota $0{,}363$,
ktorú pre model SFace odporúča knižnica OpenCV \cite{opencvzoo}; meranie na LFW ukázalo,
že pri ňom nevznikla ani jedna falošná zhoda.

\section{Návrh aktívnej výzvy}

Pasívna detekcia sa učí z~dát a jej úspešnosť v~nových podmienkach (iná kamera, iné
osvetlenie) nie je zaručená \cite{sun2023rethinking}. Preto ju dopĺňa aktívna výzva typu
\emph{challenge--response}. Server pri každom pokuse vygeneruje kryptograficky bezpečným
generátorom postupnosť
\[
  \texttt{center} \rightarrow a_1 \rightarrow a_2,\qquad
  a_1 \neq a_2 \in \{\texttt{turn\_left}, \texttt{turn\_right}, \texttt{look\_up}, \texttt{look\_down}\},
\]
teda 12 možných poradí. Výzva má 192-bitový identifikátor, platí 90~s a po prvom
použití sa zmaže. Statická fotografia nemôže vykonať pohyb a vopred nahrané video
nepozná náhodné poradie.

Poloha hlavy sa odhaduje geometricky z~piatich orientačných bodov. Body sa najprv
otočia tak, aby spojnica očí bola vodorovná (kompenzácia náklonu hlavy). Potom
\begin{equation}
  \mathrm{yaw} = \frac{x_{\mathrm{nos}} - x_{\mathrm{oči}}}{d_{\mathrm{oči}}},
  \qquad
  \mathrm{pitch} = \frac{y_{\mathrm{nos}} - y_{\mathrm{oči}}}{y_{\mathrm{ústa}} - y_{\mathrm{oči}}},
\end{equation}
kde $x_{\mathrm{oči}}, y_{\mathrm{oči}}$ je stred spojnice očí, $d_{\mathrm{oči}}$
vzdialenosť očí a $y_{\mathrm{ústa}}$ výška stredu úst. Pri otočení hlavy doľava (z
pohľadu osoby) sa nos posunie doprava v~nezrkadlenom obraze a yaw rastie; pri pohľade
hore sa nos priblíži k~línii očí a pitch klesá.

Krok \texttt{center} určí osobnú neutrálnu polohu $(\overline{\mathrm{yaw}},
\overline{\mathrm{pitch}})$ ako medián. Ostatné kroky sa hodnotia ako zmena oproti nej
-- tým sa odstráni vplyv individuálnej stavby tváre aj toho, že kamera nie je presne
pred používateľom. Krok je splnený, ak aspoň dve snímky prekročia prah ($0{,}20$ pre
yaw, $0{,}07$ pre pitch) v~správnom smere a aspoň polovica snímok kroku obsahuje
detegovanú tvár.

Prvá verzia klienta snímala každý krok pevne stanovený čas. Ak používateľ začal pohyb
neskôr, v~odoslaných snímkach ho nemuselo byť dosť a krok zlyhal, hoci ho používateľ
vykonal. Výzvu sme preto prerobili na adaptívnu: klient počas kroku posiela serveru aj
malé náhľadové snímky, server vráti aktuálnu polohu hlavy a keď pohyb prekročí prah,
klient krok o~0,6~s neskôr ukončí (najneskôr po 5~s). Toto meranie slúži iba na
navádzanie používateľa -- konečné rozhodnutie urobí server znova zo všetkých odoslaných
snímok, takže upravený klient nemôže výzvu obísť.

\section{Bezpečnostný návrh}

Tabuľka~\ref{tab:hrozby} sumarizuje uvažované hrozby a zodpovedajúce protiopatrenia.
Kategorizácia útokov vychádza z~normy ISO/IEC 30107-3 \cite{iso30107-3} a zo správy
ENISA o~vzdialenom overovaní totožnosti \cite{enisa2022remote}.

\begin{table}[H]
  \centering
  \caption{Model hrozieb a protiopatrenia}
  \label{tab:hrozby}
  \small
  \begin{tabularx}{\textwidth}{@{}>{\raggedright\arraybackslash}p{4.2cm} X@{}}
    \toprule
    Hrozba & Protiopatrenie \\
    \midrule
    Vytlačená fotografia & MiniFASNet (textúra, okraje papiera); nemožnosť vykonať pohyb \\
    Fotografia / video na displeji & MiniFASNet (moaré, rám displeja); náhodné poradie pohybov \\
    Nahrávka s~pohybmi hlavy & náhodná jednorazová výzva s~krátkou platnosťou \\
    Zámena osoby počas snímania & konzistencia vektorov príznakov všetkých snímok \\
    Opakované odoslanie požiadavky & výzva sa po použití atomicky zmaže \\
    Hrubá sila, ladenie podľa skóre & zablokovanie po 5 neúspechoch za 15~min, limit požiadaviek na IP, skóre sa v~produkcii nevracajú \\
    Zisťovanie existencie účtov & neznámy používateľ dostane rovnakú odpoveď \texttt{no\_match} \\
    Únik databázy & šifrovanie vzorov (Fernet: AES-128-CBC + HMAC-SHA256), kľúč mimo databázy \\
    Duplicitná registrácia & vyhľadanie 1:N v~galérii pri registrácii \\
    Injekcia snímok (virtuálna kamera) & iba čiastočne (PAD, výzva); úplná ochrana vyžaduje atestáciu zariadenia \\
    \bottomrule
  \end{tabularx}
\end{table}

Ochrana uložených vzorov je navrhnutá s~ohľadom na to, že biometrický údaj nemožno
po úniku zmeniť ako heslo \cite{abdullahi2024template}. Systém preto neukladá
fotografie, vzor šifruje a umožňuje používateľovi kedykoľvek vymazať účet aj vzor.
Každý pokus o~prihlásenie do účtu -- aj zamietnutý pre podvrh -- sa zobrazí v~histórii
jeho vlastníka, takže používateľ vidí, ak sa niekto pokúšal prihlásiť namiesto neho.

\chapter{Implementácia}
\label{ch:implementacia}

\section{Použité technológie}

\begin{table}[H]
  \centering
  \caption{Použité technológie}
  \label{tab:technologie}
  \begin{tabularx}{\textwidth}{@{}l X@{}}
    \toprule
    Oblasť & Technológia \\
    \midrule
    Server & Python 3.12, FastAPI, Pydantic, Uvicorn \\
    Databáza & SQLAlchemy 2, Alembic (migrácie), SQLite / PostgreSQL \\
    Biometria & OpenCV (YuNet \cite{wu2023yunet}, SFace \cite{zhong2021sface}), ONNX Runtime (MiniFASNet \cite{minivision2020silentface}) \\
    Bezpečnosť & PyJWT (prístupové tokeny), cryptography/Fernet (šifrovanie vzorov) \\
    Klient & TypeScript, React 19, React Router, Vite \\
    Kvalita kódu & pytest, Ruff, mypy, ESLint, Vitest, GitHub Actions \\
    Nasadenie & Docker, docker compose, nginx, Render (server), Vercel (klient) \\
    \bottomrule
  \end{tabularx}
\end{table}

Všetky neurónové siete sú predtrénované a voľne dostupné. Sťahuje ich skript
\path{scripts/download_models.py}, ktorý zároveň uloží kontrolné súčty SHA-256.
Spolu majú približne 42~MB a bežia na CPU.

\section{Štruktúra projektu}

\begin{lstlisting}
backend/
  app/
    api/            HTTP rozhranie (routy, závislosti)
    biometrics/     detektor, príznaky, kvalita, živosť, engine
      liveness/     pasívna detekcia, aktívna výzva, póza hlavy
    core/           konfigurácia, JWT, šifrovanie, chyby, rate-limit
    db/             ORM modely, pripojenie, migrácie
    repositories/   prístup k dátam
    services/       registrácia, overenie, identifikácia, výzvy
  evaluation/       metriky ISO/IEC, vyhodnocovacie skripty, grafy
  tests/            jednotkové a end-to-end testy
frontend/src/
  api/  auth/  biometrics/  camera/  components/  pages/
docs/               dokumentácia
report/             táto správa (LaTeX)
\end{lstlisting}

\section{Detekcia tváre a kontrola kvality}

Detektor YuNet sa volá cez triedu \texttt{cv2.FaceDetectorYN}. Pre každú tvár vráti
ohraničujúci obdĺžnik, päť orientačných bodov a skóre. Tváre sa zoradia podľa plochy
a ďalej sa pracuje s~najväčšou. Keďže objekt OpenCV si pamätá veľkosť vstupu, prístup
k~nemu chráni zámok, aby boli súbežné požiadavky bezpečné.

Kontrola kvality (trieda \texttt{QualityAssessor}) je inšpirovaná normou ISO/IEC
29794-5 \cite{iso29794-5}, ale je zjednodušená na rýchle a ľahko vysvetliteľné kritériá:
\begin{itemize}
  \item v~zábere nesmie byť ďalšia tvár väčšia ako 15~\% plochy hlavnej tváre,
  \item kratšia strana tváre aspoň 80~px,
  \item ostrosť -- rozptyl Laplaciánu na výreze tváre zmenšenom na $128\times128$~px aspoň 12,
  \item priemerný jas výrezu v~intervale $\langle 40, 225\rangle$.
\end{itemize}
Výsledkom je zoznam problémov (napr. \texttt{too\_dark}), ktoré klient zobrazí ako
krátku radu (\enquote{Je príliš tma}). Rovnaká kontrola beží aj ako samostatný endpoint
\texttt{/biometrics/quality}, ktorý klient volá počas nastavovania tváre a ktorý vracia
aj aktuálnu polohu hlavy. Prvá verzia mala hranice prísnejšie (90~px, ostrosť 25, jas
od 50). Bežné notebookové kamery však pri horšom osvetlení dávajú tmavší a menej ostrý
obraz, preto sme hranice zmiernili.

\section{Pasívna detekcia živosti}

Použité sú dva modely MiniFASNet (V2 a V1SE) z~projektu Silent-Face-Anti-Spoofing
\cite{minivision2020silentface}, exportované do formátu ONNX. Každý model dostane výrez
tváre zväčšený $2{,}7\times$ resp. $4\times$ oproti ohraničujúcemu obdĺžniku, aby videl
aj okolie tváre (rám displeja, okraj papiera), zmenšený na $80\times80$~px. Vstupom je
obraz BGR s~hodnotami 0--255 vo formáte NCHW, výstupom softmax nad troma triedami,
z~ktorých trieda~1 znamená živú tvár. Výsledné skóre je vážený priemer oboch modelov:
\begin{equation}
  p_{\mathrm{live}} = \frac{\sum_{m} w_m\, p_m(\text{live})}{\sum_m w_m}.
\end{equation}

\begin{lstlisting}[language=Python, caption={Predspracovanie a inferencia jedného modelu (zjednodušené)}]
crop = crop_with_context(image, face.bbox, scale, (80, 80))
tensor = np.transpose(crop.astype(np.float32), (2, 0, 1))[None]
logits = session.run(None, {input_name: tensor})[0]
p_live = softmax(logits)[0, LIVE_CLASS_INDEX]
\end{lstlisting}

\section{Aktívna výzva a snímanie}

Výzvy vydáva a overuje trieda \texttt{ChallengeService}. Výzva sa uloží pod náhodným
identifikátorom (\texttt{secrets.token\_urlsafe(24)}) s~dobou platnosti 90~s. Pri
overení sa výzva z~úložiska odstráni atomickou operáciou (v~Redise \texttt{GETDEL}),
takže tú istú výzvu nemožno použiť dvakrát. Overuje sa aj účel výzvy -- výzva vydaná na
prihlásenie nemôže byť použitá na registráciu.

Na strane klienta snímanie riadi hook \texttt{useLivenessCapture}. Každých 220~ms
skopíruje aktuálny snímok z~videa cez \texttt{canvas} (bez zrkadlenia, zmenšený na šírku
640~px, JPEG). Krok \texttt{center} trvá 1,6~s. V~pohybových krokoch zároveň posiela
náhľady na \texttt{/biometrics/quality} a zo vrátenej polohy hlavy počíta, koľko
z~požadovaného pohybu už používateľ urobil. Tento podiel sa zobrazuje ako vyplnenie
kruhového ukazovateľa okolo náhľadu kamery. Keď pohyb dosiahne prah, krok sa o~0,6~s
neskôr ukončí, inak najneskôr po 5~s. Z~každého kroku sa odošle najviac 10 posledných
snímok, celá relácia má teda najviac 27 snímok.

Na serveri vyhodnotenie výzvy (trieda \texttt{ActiveLivenessVerifier}) vypočíta pre
každú snímku polohu hlavy a pohyb oproti neutrálnej polohe v~smere požadovanej akcie.
Ak niektorý krok neuspeje, odpoveď obsahuje názov tohto kroku a klient používateľovi
presne povie, ktorý pohyb nerozpoznal (\enquote{Nerozpoznali sme otočenie hlavy doprava}).

\section{Registrácia, overenie a identifikácia}

Prípady použitia implementuje trieda \texttt{BiometricAuthService}. Všetky majú
rovnaký priebeh: spotrebovanie výzvy, kontrola, že snímky pokrývajú všetky kroky,
dekódovanie obrázkov s~limitom veľkosti, vyhodnotenie relácie a nakoniec akcia
podľa prípadu použitia:
\begin{itemize}
  \item \textbf{registrácia} -- kontrola obsadenosti mena, vyhľadanie tváre v~galérii
        (odmietnutie duplicitnej registrácie), uloženie používateľa a šifrovaného vzoru;
  \item \textbf{verifikácia 1:1} -- kontrola dočasného zablokovania, porovnanie
        so vzormi používateľa; neexistujúci používateľ dostane rovnakú chybu ako nezhoda;
  \item \textbf{identifikácia 1:N} -- porovnanie so všetkými vzormi aktívnych
        používateľov a výber najlepšej zhody nad prahom;
  \item \textbf{aktualizácia vzoru} -- nový vzor sa uloží, iba ak sa zhoduje
        s~pôvodným, čo bráni prevzatiu účtu po krádeži tokenu.
\end{itemize}

Každý pokus sa zapíše do auditného záznamu (režim, výsledok, dôvod zamietnutia,
skóre zhody a živosti, výsledok výzvy, počet snímok, čas spracovania, IP adresa).
Záznam sa zapisuje v~samostatnej transakcii až po potvrdení alebo zrušení hlavnej
transakcie, takže zostane zachovaný aj pri neúspechu. Z~auditu sa počíta aj
dočasné zablokovanie účtu.

\section{Bezpečnosť a ochrana údajov}

\begin{itemize}
  \item \textbf{Šifrovanie vzorov} -- vektor 128 hodnôt typu \texttt{float32} sa šifruje
        schémou Fernet. Kľúč je v~premenných prostredia, nie v~databáze. Vzor je
        viazaný na verziu modelu, po výmene modelu sa starý vzor nepoužije.
  \item \textbf{Prístupové tokeny} -- JWT podpísané HMAC-SHA256 s~platnosťou 30~min.
  \item \textbf{Ochrana proti zneužitiu} -- limit požiadaviek na IP adresu, limit veľkosti
        požiadavky (20~MB), snímky (2~MB) a počtu snímok (50); obrázky sa zmenšia na
        najviac 1280~px.
  \item \textbf{Produkčný režim} -- aplikácia odmietne štartovať bez nastavených
        tajomstiev; v~odpovediach sa nevracajú skóre, aby útočník nemohol vstup
        postupne ladiť podľa nich; vypnutá je dokumentácia API.
  \item \textbf{Klient} -- prísna politika CSP, zákaz vkladania do rámcov, token
        v~úložisku \texttt{sessionStorage} (zanikne zatvorením karty).
  \item \textbf{GDPR} -- používateľ môže vymazať účet aj vzor; pri registrácii je
        informovaný o~spracovaní biometrického údaja.
\end{itemize}

\section{Klientska aplikácia}

Pri návrhu rozhrania sme sa inšpirovali existujúcimi aplikáciami na prihlásenie tvárou,
napríklad nastavením Face ID v~systéme iOS: veľký kruhový náhľad kamery, okolo neho
prstenec z~čiarok, ktorý sa postupne vypĺňa, a pod ním jedna krátka veta s~pokynom.
Rozhranie je zámerne jednoduché -- bez vysvetľujúcich odsekov, s~jedným hlavným
tlačidlom na obrazovke. Ak je tvár v~kruhu v~poriadku, prstenec sa zafarbí nazeleno
(obr.~\ref{fig:ui-login}). Počas pohybového kroku sa vo vnútri kruhu zobrazí šípka
v~smere pohybu (obr.~\ref{fig:ui-challenge}) a po neúspechu sa prstenec zafarbí
načerveno a pod ním je dôvod zamietnutia (obr.~\ref{fig:ui-rejected}). Náhľad je
zrkadlený, aby pokyn \enquote{doľava} zodpovedal smeru, ktorý používateľ vidí.

Snímky obrazoviek v~tejto kapitole boli vytvorené s~virtuálnou kamerou, ktorej obraz
tvorila verejne dostupná, umelo vygenerovaná fotografia tváre (StyleGAN, voľné dielo).
Ide teda zároveň o~ukážku útoku statickým obrázkom -- systém ho podľa očakávania
zamietol, pretože obraz sa počas výzvy nepohol.

\begin{figure}[H]
  \centering
  \begin{subfigure}{0.49\textwidth}
    \centering
    \optfigure{ui_home.png}{\linewidth}{úvodná stránka}
    \caption{Úvodná stránka}
  \end{subfigure}
  \hfill
  \begin{subfigure}{0.49\textwidth}
    \centering
    \optfigure{ui_enroll_form.png}{\linewidth}{registrácia}
    \caption{Registrácia -- prvý krok}
  \end{subfigure}
  \caption{Úvodná stránka a registrácia}
  \label{fig:ui-home}
\end{figure}

\begin{figure}[H]
  \centering
  \optfigure{ui_login_ready.png}{0.5\textwidth}{prihlásenie, tvár v poriadku}
  \caption{Prihlásenie -- tvár je v~kruhu a spĺňa požiadavky na kvalitu}
  \label{fig:ui-login}
\end{figure}

\begin{figure}[H]
  \centering
  \begin{subfigure}{0.49\textwidth}
    \centering
    \optfigure{ui_login_challenge.png}{\linewidth}{priebeh výzvy}
    \caption{Krok výzvy \enquote{Otočte hlavu doprava}}
    \label{fig:ui-challenge}
  \end{subfigure}
  \hfill
  \begin{subfigure}{0.49\textwidth}
    \centering
    \optfigure{ui_login_rejected.png}{\linewidth}{zamietnutie}
    \caption{Zamietnutie statického obrázka}
    \label{fig:ui-rejected}
  \end{subfigure}
  \caption{Aktívna výzva počas prihlásenia}
\end{figure}

\begin{figure}[H]
  \centering
  \begin{subfigure}{0.49\textwidth}
    \centering
    \optfigure{ui_dashboard.png}{\linewidth}{účet}
    \caption{Účet s~históriou pokusov}
  \end{subfigure}
  \hfill
  \begin{subfigure}{0.49\textwidth}
    \centering
    \optfigure{ui_admin.png}{\linewidth}{správa}
    \caption{Správa systému}
  \end{subfigure}
  \caption{Používateľský účet a administrácia}
  \label{fig:ui-account}
\end{figure}

Technicky sa klient skladá z~troch častí. Hook \texttt{useCamera} spúšťa a zastavuje
webkameru a prekladá chyby prehliadača na zrozumiteľné hlásenia. Pri testovaní sme
v~ňom našli súbeh: ak sa kamera spustila dvakrát rýchlo po sebe (napríklad pri rýchlom
prechode medzi stránkami), druhé spustenie prerušilo prvé a aplikácia nesprávne
hlásila, že kamera nie je dostupná. Oprava každé spustenie očísluje a výsledky
neaktuálnych spustení zahodí. Hook \texttt{useQualityMonitor} pred začiatkom snímania
pravidelne kontroluje kvalitu a \texttt{useLivenessCapture} riadi samotnú výzvu.
Chyby zo servera sa podľa kódu a dôvodu zamietnutia prekladajú na krátke slovenské
hlásenia.

\section{Nasadenie}

Server aj klient majú vlastný \texttt{Dockerfile}. Lokálne sa celý systém spustí
príkazom \texttt{docker compose up}: kontajner nginx poskytuje statické súbory klienta
a presmerúva \texttt{/api} na server, server pri štarte spustí databázové migrácie
a Redis zdieľa výzvy medzi procesmi. Verejná inštancia beží na službách Render
(server s~PostgreSQL, definícia v~\texttt{render.yaml}) a Vercel (klient,
\texttt{vercel.json}), pričom Vercel presmerúva \texttt{/api} na server -- prehliadač
tak komunikuje iba s~jednou doménou.

\begin{itemize}
  \item Aplikácia: \url{https://face-auth-antispoofing.vercel.app}
  \item Repozitár: \url{https://github.com/opoiasnik/face-auth-antispoofing}
\end{itemize}

\chapter{Testovanie a vyhodnotenie}
\label{ch:testovanie}

Systém sme overili na štyroch úrovniach: (1)~automatické testy softvéru, (2)~meranie
presnosti rozpoznávania na verejnej databáze LFW, (3)~meranie pasívnej detekcie útokov
na verejnej databáze CASIA-FASD a (4)~testy celého systému cez jeho rozhranie. Výkon
sme merali samostatne. Všetky vyhodnocovacie skripty sú v~priečinku
\texttt{backend/evaluation} a výsledky ukladajú ako tabuľky (JSON, CSV, Markdown)
a grafy (PNG, PDF).

\section{Metriky}

\subsection*{Rozpoznávanie tváre}
Pri prahu $\tau$ a skóre porovnaní genuine $G$ (rovnaká osoba) a impostor $I$ (rôzne
osoby):
\begin{equation}
  \mathrm{FAR}(\tau) = \frac{|\{s \in I : s \ge \tau\}|}{|I|},
  \qquad
  \mathrm{FRR}(\tau) = \frac{|\{s \in G : s < \tau\}|}{|G|}.
\end{equation}
Miera rovnakej chyby (EER) je hodnota v~bode $\mathrm{FAR}(\tau) = \mathrm{FRR}(\tau)$;
počíta sa lineárnou interpoláciou medzi susednými prahmi. Ďalej uvádzame FRR pri
pevnom FAR (1~\% a 0,1~\%) a podiel snímok bez detegovanej tváre (FTA).

\subsection*{Detekcia prezentačných útokov}
Podľa ISO/IEC 30107-3 \cite{iso30107-3} sa pre každý druh útoku $k$ počíta
\begin{equation}
  \mathrm{APCER}_k = \frac{\text{útoky typu } k \text{ prijaté ako živá tvár}}{\text{všetky útoky typu } k},
  \qquad
  \mathrm{BPCER} = \frac{\text{odmietnuté živé tváre}}{\text{všetky živé tváre}},
\end{equation}
pričom celkové $\mathrm{APCER} = \max_k \mathrm{APCER}_k$. Uvádzame aj BPCER20
a BPCER100, teda BPCER pri prahu, pri ktorom $\mathrm{APCER} \le 5~\%$ resp. $1~\%$,
a pre porovnanie s~literatúrou aj ACER.

\section{Automatické testy}

Serverová časť má 49 testov (pytest) s~pokrytím kódu 87~\%. Testy nahrádzajú neurónové
siete deterministickými náhradami: syntetická snímka nesie v~jednom pixeli identitu,
polohu hlavy a príznak živosti. Vďaka tomu sa celý tok overuje cez skutočné HTTP
rozhranie bez modelov a bez GPU. Klient má 7 jednotkových testov v~nástroji Vitest.
Testy, statická analýza (Ruff, mypy, ESLint) a zostavenie Docker obrazov bežia
automaticky pri každej zmene v~GitHub Actions.

\begin{table}[H]
  \centering
  \caption{Pokrytie automatickými testami}
  \label{tab:testy}
  \small
  \begin{tabularx}{\textwidth}{@{}l X@{}}
    \toprule
    Oblasť & Overované správanie \\
    \midrule
    Poloha hlavy & znamienko yaw a pitch, nezávislosť od náklonu hlavy, degenerované body, výpočet polohy aj z~rozmazanej snímky \\
    Aktívna výzva & správny a nesprávny smer, nedostatočný pohyb, málo snímok s~tvárou, náhodnosť poradia \\
    Metriky & FAR/FRR, EER (aj analyticky pre dve normálne rozdelenia), APCER/BPCER/ACER \\
    Bezpečnosť & podpis a platnosť JWT, šifrovanie vzorov, produkčná konfigurácia, rate-limit \\
    Výzvy & jednorazovosť, expirácia, kontrola účelu \\
    End-to-end API & registrácia a overenie, impostor, neznámy používateľ, útok fotografiou (pasívny aj aktívny), zámena osoby počas relácie, opakované použitie výzvy, chýbajúci krok, duplicitná tvár, identifikácia 1:N, zablokovanie účtu, výmaz účtu, práva administrátora, história zamietnutých pokusov \\
    Klient & výpočet pohybu hlavy, medián neutrálnej polohy, preklad chybových hlásení \\
    \bottomrule
  \end{tabularx}
\end{table}

\section{Presnosť rozpoznávania tváre}

Model SFace sme vyhodnotili na databáze Labeled Faces in the Wild \cite{huang2007lfw}
podľa protokolu \emph{pairs} -- 6000 dvojíc fotografií, z~toho 3000 dvojíc tej istej
osoby a 3000 dvojíc rôznych osôb. Kontrolu kvality sme pri tomto meraní vypli, aby sme
merali samotný model rozpoznávania. Tvár sa našla na všetkých 12\,000 fotografiách
(FTA 0~\%). Výsledky sú v~tabuľke~\ref{tab:lfw} a na obr.~\ref{fig:lfw}.

\begin{table}[H]
  \centering
  \caption{Výsledky rozpoznávania tváre na LFW (6000 dvojíc)}
  \label{tab:lfw}
  \begin{tabular}{@{}lccc@{}}
    \toprule
    Pracovný bod & Prah & FAR [\%] & FRR [\%] \\
    \midrule
    EER & 0,266 & 2,47 & 2,47 \\
    Systémový prah & 0,400 & 0,00 & 3,90 \\
    FAR $\le$ 1~\% & 0,294 & 1,00 & 2,63 \\
    FAR $\le$ 0,1~\% & 0,339 & 0,10 & 3,03 \\
    \bottomrule
  \end{tabular}
\end{table}

\begin{figure}[H]
  \centering
  \begin{subfigure}{0.49\textwidth}
    \centering
    \optfigure{score_distribution.pdf}{\linewidth}{histogram skóre}
    \caption{Rozdelenie skóre}
  \end{subfigure}
  \hfill
  \begin{subfigure}{0.49\textwidth}
    \centering
    \optfigure{det.pdf}{\linewidth}{DET krivka}
    \caption{DET krivka}
  \end{subfigure}
  \caption{Rozpoznávanie tváre na LFW}
  \label{fig:lfw}
\end{figure}

Rozdelenia skóre dvojíc tej istej osoby a rôznych osôb sa prekrývajú len málo. Pri
použitom prahu 0,40 nevznikla v~3000 porovnaniach rôznych osôb ani jedna falošná zhoda,
cenou je odmietnutie 3,9~\% dvojíc tej istej osoby. Databáza LFW obsahuje fotografie
z~rôznych rokov, s~rôznou pózou a osvetlením, čo je náročnejšia situácia ako
prihlasovanie pred vlastnou webkamerou. Pri prihlasovaní sa navyše vzor relácie počíta z~viacerých
snímok, čo náhodné výkyvy ďalej zmierňuje. Prah 0,40 preto považujeme za vhodný
kompromis medzi bezpečnosťou a pohodlím.

\section{Presnosť pasívnej detekcie útokov}

Na vyhodnotenie sme použili testovaciu časť verejnej databázy CASIA-FASD
\cite{zhang2012casia}, konkrétne 2408 snímok vybraných z~jej videí. Databáza obsahuje
živé tváre a tri druhy útokov (tab.~\ref{tab:pad-dataset}). Snímky v~použitej verzii
databázy sú už orezané okolo tváre na $256\times256$~px, takže modely MiniFASNet nevidia
celé okolie tváre, s~ktorým počítajú. Výsledky sú preto skôr pesimistickým odhadom.

\begin{table}[H]
  \centering
  \caption{Dataset pre vyhodnotenie pasívnej detekcie (CASIA-FASD, testovacia časť)}
  \label{tab:pad-dataset}
  \begin{tabularx}{\textwidth}{@{}l X r@{}}
    \toprule
    Trieda & Popis & Snímky \\
    \midrule
    \texttt{bona\_fide} & živé tváre & 591 \\
    \texttt{print\_warped} & vytlačená fotografia, ohýbaná pred kamerou & 712 \\
    \texttt{print\_cut} & vytlačená fotografia s~vystrihnutými očami & 536 \\
    \texttt{replay} & video prehrávané na displeji & 569 \\
    \bottomrule
  \end{tabularx}
\end{table}

Pri pôvodnom prahu 0,70 boli útoky zachytené takmer úplne (APCER najviac 0,93~\%), no
odmietnutých bolo 34~\% snímok živých tvárí. Taký prah by v~praxi viedol k~častému
odmietaniu skutočných používateľov. Tabuľka
\ref{tab:prah} ukazuje, ako sa chyby menia s~prahom.

\begin{table}[H]
  \centering
  \caption{Chyby pasívnej detekcie pri rôznych prahoch (úroveň snímok)}
  \label{tab:prah}
  \begin{tabular}{@{}ccccc@{}}
    \toprule
    Prah & BPCER [\%] & APCER print\_warped & APCER print\_cut & APCER replay \\
    \midrule
    0,20 & 5,58 & 3,93 & 5,97 & 5,62 \\
    0,30 & 8,63 & 2,67 & 2,99 & 2,64 \\
    0,40 & 12,86 & 1,83 & 2,05 & 0,88 \\
    0,50 & 18,10 & 0,70 & 1,68 & 0,53 \\
    0,60 & 25,21 & 0,28 & 1,12 & 0,18 \\
    0,70 & 34,35 & 0,00 & 0,93 & 0,00 \\
    \bottomrule
  \end{tabular}
\end{table}

Prah sme znížili na 0,50. Na úrovni jednej snímky pri ňom APCER neprekročí 1,7~\%
pri žiadnom type útoku. BPCER 18~\% sa týka jednotlivých snímok -- pri prihlásení sa
skóre priemeruje cez približne sedem frontálnych snímok, takže náhodne nízke skóre
jednej snímky reláciu nezamietne. Útoky, ktoré by pasívna detekcia prepustila, musia
navyše prejsť aktívnou výzvou. Výsledky pri systémovom prahu sú v~tabuľke~\ref{tab:pad}.

\begin{table}[H]
  \centering
  \caption{Výsledky pasívnej detekcie útokov na CASIA-FASD (prah 0,50)}
  \label{tab:pad}
  \begin{tabular}{@{}lc@{}}
    \toprule
    Metrika & Hodnota [\%] \\
    \midrule
    APCER (print\_warped) & 0,70 \\
    APCER (print\_cut) & 1,68 \\
    APCER (replay) & 0,53 \\
    APCER (max) & 1,68 \\
    BPCER & 18,10 \\
    ACER & 9,89 \\
    EER & 5,25 \\
    BPCER20 (APCER = 5~\%) & 5,58 \\
    BPCER100 (APCER = 1~\%) & 16,58 \\
    \bottomrule
  \end{tabular}
\end{table}

\begin{figure}[H]
  \centering
  \begin{subfigure}{0.49\textwidth}
    \centering
    \optfigure{pad_score_distribution.pdf}{\linewidth}{skóre živosti podľa triedy}
    \caption{Rozdelenie skóre živosti}
  \end{subfigure}
  \hfill
  \begin{subfigure}{0.49\textwidth}
    \centering
    \optfigure{pad_det.pdf}{\linewidth}{DET pre každý typ útoku}
    \caption{DET (APCER vs. BPCER)}
  \end{subfigure}
  \caption{Pasívna detekcia prezentačných útokov na CASIA-FASD}
  \label{fig:pad}
\end{figure}

Z~histogramu je vidieť, že skóre takmer všetkých útokov je blízko nuly a útoky sa líšia
len v~tom, ako často sa objaví výnimka. Najťažšie odhaliteľná je fotografia s~vystrihnutými
očami, pri ktorej sú v~obraze skutočné oči útočníka. Živé tváre majú skóre rozložené
širšie, čo súvisí s~nižšou kvalitou videí v~databáze a s~orezaním okolia tváre.

\section{Testy celého systému}

\subsection*{Útok statickým obrázkom}
Útok statickým obrázkom sme vyskúšali dvoma spôsobmi. V~prvom sme do prehliadača
podstrčili virtuálnu kameru, ktorej obraz tvorila umelo vygenerovaná fotografia tváre,
a prešli sme bežným prihlásením cez webové rozhranie. Pasívna detekcia v~tomto prípade
obraz nezamietla (skóre živosti 0,98) -- obrázok vložený priamo do video toku nemá
žiadne stopy tlače ani displeja, ide teda v~skutočnosti o~injekčný útok. Zachytila ho
až aktívna výzva, pretože tvár sa v~žiadnom kroku nepohla
(obr.~\ref{fig:ui-rejected}). Počas tvorby snímok obrazovky sme takto vykonali 6 pokusov
a všetky boli zamietnuté s~dôvodom \texttt{active\_challenge\_failed}.

V~druhom teste sme ten istý obrázok poslali 30-krát priamo cez API (režim identifikácie),
s~drobným posunom a šumom medzi snímkami a so stratovou kompresiou JPEG podobne ako pri
skutočnej kamere. Všetkých 30 relácií bolo zamietnutých už pasívnou detekciou
(\texttt{passive\_spoof}); priemerné skóre živosti relácie bolo 0,25 (od 0,04 do 0,52).
Oba testy ukazujú, prečo systém používa obe vrstvy: každá zachytila prípad, ktorý by
druhá prepustila.

\begin{table}[H]
  \centering
  \caption{Útok statickým obrázkom}
  \label{tab:static}
  \begin{tabularx}{\textwidth}{@{}X c c l@{}}
    \toprule
    Spôsob útoku & Pokusy & Zamietnuté & Vrstva, ktorá útok zachytila \\
    \midrule
    Virtuálna kamera v~prehliadači & 6 & 6 & aktívna výzva \\
    Priame volanie API & 30 & 30 & pasívna detekcia \\
    \bottomrule
  \end{tabularx}
\end{table}

\subsection*{Test so živými používateľmi}
Úspešnosť prihlásenia skutočných používateľov a odolnosť proti fyzickým útokom
(vytlačená fotografia, displej mobilu, video s~pohybmi hlavy) sa dá zmerať iba
s~reálnymi ľuďmi pred kamerou. Postup a scenáre sú v~tabuľke~\ref{tab:e2e}; výsledky sa
čerpajú z~auditného záznamu skriptom \texttt{python -m evaluation.audit\_report}.

\begin{table}[H]
  \centering
  \caption{Scenáre testu so živými používateľmi}
  \label{tab:e2e}
  \small
  \begin{tabularx}{\textwidth}{@{}X c l c@{}}
    \toprule
    Scenár & Pokusy & Očakávaný výsledok & Úspešnosť [\%] \\
    \midrule
    Registrovaný používateľ, 1:1 & \todo{} & prijatie & \todo{} \\
    Registrovaný používateľ, 1:N & \todo{} & prijatie & \todo{} \\
    Iná osoba s~menom obete & \todo{} & \texttt{no\_match} & \todo{} \\
    Vytlačená fotografia obete & \todo{} & zamietnutie & \todo{} \\
    Fotografia na mobile/monitore & \todo{} & zamietnutie & \todo{} \\
    Video obete s~pohybmi hlavy & \todo{} & zamietnutie & \todo{} \\
    Zámena osoby počas snímania & \todo{} & \texttt{identity\_inconsistent} & \todo{} \\
    \bottomrule
  \end{tabularx}
\end{table}

\section{Výkon}

Latenciu sme merali skriptom \texttt{evaluation.benchmark} na notebooku s~procesorom
Intel Core i7-13650HX (ONNX Runtime na CPU, bez GPU), 200 opakovaní, na snímke
v~rozlíšení $640\times360$~px, ktoré posiela klient.

\begin{table}[H]
  \centering
  \caption{Latencia spracovania jednej snímky ($640\times360$~px)}
  \label{tab:latencia}
  \begin{tabular}{@{}lrrr@{}}
    \toprule
    Fáza & Priemer [ms] & Medián [ms] & P95 [ms] \\
    \midrule
    Detekcia (YuNet) & 4,76 & 4,65 & 5,50 \\
    Kontrola kvality & 0,10 & 0,10 & 0,11 \\
    Odhad polohy hlavy & 0,01 & 0,01 & 0,01 \\
    Pasívna živosť (2$\times$ MiniFASNet) & 2,59 & 2,53 & 2,90 \\
    Extrakcia príznakov (SFace) & 3,98 & 3,95 & 4,38 \\
    \midrule
    Celá snímka & 12,23 & 12,06 & 13,31 \\
    \bottomrule
  \end{tabular}
\end{table}

Relácia má najviac 27 snímok, jej spracovanie na serveri teda trvá približne 0,3~s.
Pre porovnanie, pri snímke $1280\times720$~px trvá samotná detekcia 18~ms a celá snímka
26~ms, preto klient snímky pred odoslaním zmenšuje. Celkovú dobu prihlásenia (približne
5--7~s) určuje najmä čas, ktorý používateľ potrebuje na vykonanie pohybov. Na
bezplatnej inštancii služby Render je spracovanie pomalšie a prvá požiadavka po
nečinnosti trvá až desiatky sekúnd, kým sa inštancia spustí.

\chapter{Príručka a dokumentácia}
\label{ch:prirucka}

\section{Archivácia a reprodukovateľnosť}

Zdrojový kód, dokumentácia, vyhodnocovacie skripty aj táto správa sú v~repozitári
\url{https://github.com/opoiasnik/face-auth-antispoofing}. Reprodukovateľnosť
zabezpečujú:
\begin{itemize}
  \item skript na stiahnutie modelov s~kontrolnými súčtami SHA-256,
  \item databázové migrácie (Alembic) -- schéma sa vytvorí automaticky,
  \item uzamknuté verzie závislostí klienta (\texttt{package-lock.json}),
  \item Docker obrazy a súbor \texttt{docker-compose.yml} pre spustenie jedným príkazom,
  \item automatické testy spúšťané v~GitHub Actions pri každej zmene,
  \item vyhodnocovacie skripty, ktoré sami stiahnu verejné dáta (LFW) alebo pracujú
        s~priečinkom s~jasne danou štruktúrou (CASIA-FASD, vlastné nahrávky).
\end{itemize}

Podrobná technická dokumentácia je v~priečinku \texttt{docs/} repozitára:
architektúra (\texttt{architecture.md}), popis REST API (\texttt{api.md}),
používateľská príručka (\texttt{user-guide.md}) a metodika vyhodnotenia
(\texttt{evaluation.md}).

\section{Inštalácia a spustenie}

\subsection*{Lokálne spustenie}
Požiadavky: Python 3.11 alebo novší, Node.js 20 alebo novší, webkamera.

\begin{lstlisting}[language=bash, caption={Spustenie servera}]
cd backend
python -m venv .venv
.venv/Scripts/activate        # Linux/macOS: source .venv/bin/activate
pip install -e ".[dev,eval]"
python scripts/download_models.py
uvicorn app.main:create_app --factory --port 8000
\end{lstlisting}

\begin{lstlisting}[language=bash, caption={Spustenie klienta}]
cd frontend
npm ci
npm run dev                   # http://localhost:5173
\end{lstlisting}

\subsection*{Spustenie v~Dockeri}
\begin{lstlisting}[language=bash]
cp .env.example .env
docker compose up -d --build  # http://localhost:8080
\end{lstlisting}

Pri produkčnom nasadení treba nastaviť \texttt{FA\_ENVIRONMENT=production},
vygenerovať tajomstvá príkazom \texttt{python -m app.cli generate-secrets} a aplikáciu
sprístupniť cez HTTPS -- prehliadač inak nepovolí prístup ku kamere.

\subsection*{Vyhodnotenie}
\begin{lstlisting}[language=bash]
python -m evaluation.evaluate_recognition --lfw
python -m evaluation.evaluate_pad --dir data/pad
python -m evaluation.benchmark --image tvar.jpg --runs 200
python -m evaluation.audit_report
\end{lstlisting}

\subsection*{Konfigurácia}
Prahy a limity sa nastavujú premennými prostredia s~predponou \texttt{FA\_}
(tab.~\ref{tab:konfiguracia}).

\begin{table}[H]
  \centering
  \caption{Najdôležitejšie konfiguračné premenné}
  \label{tab:konfiguracia}
  \small
  \begin{tabularx}{\textwidth}{@{}l c X@{}}
    \toprule
    Premenná & Predvolená & Význam \\
    \midrule
    \texttt{FA\_MATCHING\_\_THRESHOLD} & 0,40 & prah kosínusovej podobnosti \\
    \texttt{FA\_PASSIVE\_LIVENESS\_\_THRESHOLD} & 0,50 & prah skóre živosti \\
    \texttt{FA\_ACTIVE\_LIVENESS\_\_STEPS} & 2 & počet náhodných pohybov vo výzve \\
    \texttt{FA\_ACTIVE\_LIVENESS\_\_YAW\_THRESHOLD} & 0,20 & minimálne otočenie hlavy \\
    \texttt{FA\_ACTIVE\_LIVENESS\_\_PITCH\_THRESHOLD} & 0,07 & minimálny pohyb hore/dole \\
    \texttt{FA\_CHALLENGE\_\_ACTION\_TIMEOUT\_MS} & 5000 & najdlhšie trvanie pohybového kroku \\
    \texttt{FA\_SECURITY\_\_LOCKOUT\_MAX\_FAILURES} & 5 & neúspechy do zablokovania \\
    \texttt{FA\_SECURITY\_\_EXPOSE\_SCORES} & true & vracanie skóre (v~produkcii false) \\
    \texttt{FA\_DATABASE\_URL} & SQLite & pripojenie k~databáze \\
    \texttt{FA\_REDIS\_URL} & -- & zdieľané výzvy pri viacerých procesoch \\
    \bottomrule
  \end{tabularx}
\end{table}

\section{Používateľská príručka}

\subsection*{Registrácia}
\begin{enumerate}
  \item Na stránke \emph{Registrovať} zadajte používateľské meno (3--32 znakov: malé
        písmená, číslice, bodka, pomlčka, podčiarkovník), voliteľne meno a kliknite na
        \emph{Pokračovať}.
  \item Povoľte prístup ku kamere a umiestnite tvár do kruhu. Keď sa prstenec okolo kruhu
        zafarbí nazeleno a zobrazí sa \enquote{Tvár je v~poriadku}, kliknite na
        \emph{Začať}.
  \item Pozrite sa do kamery a potom urobte pohyby podľa pokynu pod kruhom (napríklad
        \enquote{Otočte hlavu doľava}); smer ukazuje aj šípka v~kruhu. Krok sa ukončí
        sám, keď systém pohyb rozpozná. Poradie pohybov je pri každom pokuse iné.
  \item Po úspešnom overení sa vytvorí účet a používateľ je prihlásený.
\end{enumerate}

\subsection*{Prihlásenie}
Štandardne sa zadáva meno (verifikácia 1:1). Odkaz \emph{Prihlásiť sa bez mena}
prepne na identifikáciu 1:N iba podľa tváre. Postup snímania je rovnaký ako pri
registrácii. Po piatich neúspešných pokusoch počas 15 minút sa účet dočasne zablokuje.

\subsection*{Riešenie problémov}
\begin{table}[H]
  \centering
  \caption{Časté hlásenia a ich riešenie}
  \label{tab:hlasenia}
  \small
  \begin{tabularx}{\textwidth}{@{}>{\raggedright\arraybackslash}p{5.5cm} X@{}}
    \toprule
    Hlásenie & Riešenie \\
    \midrule
    Kamera nie je dostupná & povoľte kameru v~nastaveniach stránky a kliknite na \emph{Povoliť kameru} \\
    Nevidíme vašu tvár / Podíďte bližšie & tvár do kruhu, bližšie ku kamere \\
    Je príliš tma / Je príliš svetla & otočte sa k~zdroju svetla, nesedte proti oknu \\
    Nerozpoznali sme \ldots{} & urobte uvedený pohyb výraznejšie, až keď sa zobrazí pokyn \\
    Vyzerá to na fotografiu alebo obrazovku & nepoužívajte fotografiu ani displej; pri falošnom poplachu zlepšite osvetlenie \\
    Tvár sa nezhoduje s~účtom & zlepšite svetlo; pri zmene vzhľadu aktualizujte vzor v~účte \\
    Príliš veľa pokusov & počkajte 15 minút \\
    \bottomrule
  \end{tabularx}
\end{table}

\subsection*{Používateľský účet a administrácia}
V~účte je história všetkých pokusov o~prihlásenie vrátane neúspešných, možnosť
aktualizovať vzor tváre (nový vzor sa musí zhodovať s~pôvodným) a natrvalo vymazať
účet aj vzor (obr.~\ref{fig:ui-account}). Rolu administrátora prideľuje príkaz
\texttt{python -m app.cli promote <meno>}. Administrátor vidí štatistiky úspešnosti,
dôvody zamietnutí, zoznam používateľov (blokovanie, odstránenie) a celý auditný záznam.

\section{Rozhranie REST API}

\begin{table}[H]
  \centering
  \caption{Hlavné koncové body API (predpona \texttt{/api/v1})}
  \label{tab:api}
  \small
  \begin{tabularx}{\textwidth}{@{}l l X@{}}
    \toprule
    Metóda & Cesta & Popis \\
    \midrule
    POST & \texttt{/challenges} & vydanie jednorazovej výzvy \\
    POST & \texttt{/biometrics/quality} & kvalita a poloha hlavy na jednej snímke \\
    POST & \texttt{/enrollment} & registrácia, vráti JWT \\
    POST & \texttt{/auth/verify} & overenie 1:1, vráti JWT \\
    POST & \texttt{/auth/identify} & identifikácia 1:N, vráti JWT \\
    GET & \texttt{/users/me/attempts} & história vlastných pokusov \\
    PUT & \texttt{/users/me/face} & aktualizácia vzoru \\
    DELETE & \texttt{/users/me} & výmaz účtu a vzoru \\
    GET & \texttt{/admin/stats}, \texttt{/admin/attempts} & štatistiky a audit (administrátor) \\
    \bottomrule
  \end{tabularx}
\end{table}

Úplný popis vrátane formátu chýb je v~súbore \texttt{docs/api.md} a v~interaktívnej
dokumentácii \texttt{/docs} (dostupná mimo produkčného režimu).

\chapter{Záver}

V~projekte sme navrhli, implementovali a nasadili webový systém na prihlasovanie tvárou
s~viacvrstvovou ochranou proti prezentačným útokom. Systém tvár sníma bežnou webkamerou,
všetky rozhodnutia robí na serveri a kombinuje kontrolu kvality snímky, pasívnu detekciu
podvrhu modelmi MiniFASNet, náhodnú výzvu na pohyb hlavy a kontrolu konzistencie
identity počas snímania. Tvár porovnáva modelom SFace a uložené vzory šifruje.

Meranie ukázalo, že rozpoznávanie dosahuje na databáze LFW EER 2,47~\% a pri použitom
prahu 0,40 nevzniklo ani jedno falošné prijatie. Pasívna detekcia mala na databáze
CASIA-FASD EER 5,25~\%. Meranie zároveň odhalilo, že pôvodný prah 0,70 zamietal príliš
veľa živých tvárí, preto sme ho znížili na 0,50, pri ktorom APCER neprekročí 1,7~\%.
Pri teste statickým obrázkom sa ukázalo, že vrstvy sa dopĺňajú: obraz vložený cez
virtuálnu kameru prešiel pasívnou detekciou, ale zastavila ho aktívna výzva, a naopak
snímky poslané priamo cez API zachytila pasívna detekcia. Spracovanie relácie trvá na
bežnom procesore približne 0,3~s.

Počas testovania sme na základe nameraných výsledkov upravili aj samotný systém:
výzva sa namiesto pevného času prispôsobuje pohybu používateľa, poloha hlavy sa počíta
aj z~rozmazaných snímok, pasívna detekcia sa vyhodnocuje len na frontálnych snímkach
a opravili sme chybu pri spúšťaní kamery.

Systém má aj obmedzenia. Poloha hlavy sa odhaduje len z~piatich bodov, takže je menej
presná ako odhad z~hustej siete bodov. Proti injekčným útokom cez virtuálnu kameru
chráni iba aktívna výzva -- dostatočne sofistikovaný útok, ktorý by v~reálnom čase
generoval pohyby tváre, by vyžadoval ďalšie opatrenia, napríklad atestáciu zariadenia.
Pasívnu detekciu sme vyhodnotili na staršej databáze s~orezanými snímkami a test so
živými používateľmi v~širšom rozsahu zostáva na doplnenie. Ako rozšírenie sa ponúka
použitie modelov detekcie s~lepšou generalizáciou \cite{srivatsan2023flip}, detekcia
injekčných útokov a ochrana vzorov zrušiteľnou biometriou \cite{abdullahi2024template}.

\begin{thebibliography}{99}
\addcontentsline{toc}{chapter}{Zoznam použitej literatúry}

\bibitem{enisa2022remote}
EUROPEAN UNION AGENCY FOR CYBERSECURITY (ENISA).
\textit{Remote Identity Proofing: Attacks \& Countermeasures}.
Athens: ENISA, 2022.
Dostupné z: \url{https://www.enisa.europa.eu/publications/remote-identity-proofing-attacks-countermeasures}

\bibitem{wu2023yunet}
WU, Wei, PENG, Hanyang a YU, Shiqi.
YuNet: A Tiny Millisecond-level Face Detector.
\textit{Machine Intelligence Research}. 2023, roč.~20, č.~5, s.~656--665.
DOI: 10.1007/s11633-023-1423-y

\bibitem{deng2022arcface}
DENG, Jiankang, GUO, Jia, YANG, Jing, XUE, Niannan, KOTSIA, Irene a ZAFEIRIOU, Stefanos.
ArcFace: Additive Angular Margin Loss for Deep Face Recognition.
\textit{IEEE Transactions on Pattern Analysis and Machine Intelligence}. 2022,
roč.~44, č.~10, s.~5962--5979.
DOI: 10.1109/TPAMI.2021.3087709

\bibitem{kim2022adaface}
KIM, Minchul, JAIN, Anil K. a LIU, Xiaoming.
AdaFace: Quality Adaptive Margin for Face Recognition.
In: \textit{2022 IEEE/CVF Conference on Computer Vision and Pattern Recognition (CVPR)}.
IEEE, 2022, s.~18729--18738.
DOI: 10.1109/CVPR52688.2022.01819

\bibitem{zhong2021sface}
ZHONG, Yaoyao, DENG, Weihong, HU, Jiani, ZHAO, Dongyue, LI, Xian a WEN, Dongchao.
SFace: Sigmoid-Constrained Hypersphere Loss for Robust Face Recognition.
\textit{IEEE Transactions on Image Processing}. 2021, roč.~30, s.~2587--2598.
DOI: 10.1109/TIP.2020.3048632

\bibitem{schlett2022fiqa}
SCHLETT, Torsten, RATHGEB, Christian, HENNIGER, Olaf, GALBALLY, Javier,
FIERREZ, Julian a BUSCH, Christoph.
Face Image Quality Assessment: A Literature Survey.
\textit{ACM Computing Surveys}. 2022, roč.~54, č.~10s, s.~1--49.
DOI: 10.1145/3507901

\bibitem{iso29794-5}
ISO/IEC 29794-5:2025.
\textit{Information technology -- Biometric sample quality -- Part~5: Face image data}.
1.~vyd. Ženeva: ISO, 2025.
Dostupné z: \url{https://www.iso.org/standard/81005.html}

\bibitem{hernandezortega2023pad}
HERNANDEZ-ORTEGA, Javier, FIERREZ, Julian, MORALES, Aythami a GALBALLY, Javier.
Introduction to Presentation Attack Detection in Face Biometrics and Recent Advances.
In: MARCEL, Sébastien, FIERREZ, Julian a EVANS, Nicholas (ed.).
\textit{Handbook of Biometric Anti-Spoofing: Presentation Attack Detection and
Vulnerability Assessment}. 3.~vyd. Singapore: Springer, 2023, s.~203--230.
DOI: 10.1007/978-981-19-5288-3\_9

\bibitem{yu2023fassurvey}
YU, Zitong, QIN, Yunxiao, LI, Xiaobai, ZHAO, Chenxu, LEI, Zhen a ZHAO, Guoying.
Deep Learning for Face Anti-Spoofing: A Survey.
\textit{IEEE Transactions on Pattern Analysis and Machine Intelligence}. 2023,
roč.~45, č.~5, s.~5609--5631.
DOI: 10.1109/TPAMI.2022.3215850

\bibitem{sun2023rethinking}
SUN, Yiyou, LIU, Yaojie, LIU, Xiaoming, LI, Yixuan a CHU, Wen-Sheng.
Rethinking Domain Generalization for Face Anti-spoofing: Separability and Alignment.
In: \textit{2023 IEEE/CVF Conference on Computer Vision and Pattern Recognition (CVPR)}.
IEEE, 2023, s.~24563--24574.
DOI: 10.1109/CVPR52729.2023.02353

\bibitem{srivatsan2023flip}
SRIVATSAN, Koushik, NASEER, Muzammal a NANDAKUMAR, Karthik.
FLIP: Cross-domain Face Anti-spoofing with Language Guidance.
In: \textit{2023 IEEE/CVF International Conference on Computer Vision (ICCV)}.
IEEE, 2023, s.~19628--19639.
DOI: 10.1109/ICCV51070.2023.01803

\bibitem{wang2023wildfas}
WANG, Dong, GUO, Jia, SHAO, Qiqi et al.
Wild Face Anti-Spoofing Challenge 2023: Benchmark and Results.
In: \textit{2023 IEEE/CVF Conference on Computer Vision and Pattern Recognition
Workshops (CVPRW)}. IEEE, 2023, s.~6380--6391.
DOI: 10.1109/CVPRW59228.2023.00679

\bibitem{iso30107-3}
ISO/IEC 30107-3:2023.
\textit{Information technology -- Biometric presentation attack detection --
Part~3: Testing and reporting}. 2.~vyd. Ženeva: ISO, 2023.
Dostupné z: \url{https://www.iso.org/standard/79520.html}

\bibitem{busch2023standards}
BUSCH, Christoph.
Standards for Biometric Presentation Attack Detection.
In: MARCEL, Sébastien, FIERREZ, Julian a EVANS, Nicholas (ed.).
\textit{Handbook of Biometric Anti-Spoofing: Presentation Attack Detection and
Vulnerability Assessment}. 3.~vyd. Singapore: Springer, 2023, s.~571--583.
DOI: 10.1007/978-981-19-5288-3\_21

\bibitem{abdullahi2024template}
ABDULLAHI, Sani M., SUN, Shuifa, WANG, Beng, WEI, Ning a WANG, Hongxia.
Biometric Template Attacks and Recent Protection Mechanisms: A Survey.
\textit{Information Fusion}. 2024, roč.~103, čl.~102144.
DOI: 10.1016/j.inffus.2023.102144

\bibitem{euaiact2024}
EURÓPSKY PARLAMENT A RADA EURÓPSKEJ ÚNIE.
Nariadenie (EÚ) 2024/1689 z~13.~júna 2024, ktorým sa stanovujú harmonizované pravidlá
v~oblasti umelej inteligencie (akt o~umelej inteligencii).
\textit{Úradný vestník Európskej únie}, L, 2024/1689.
Dostupné z: \url{https://eur-lex.europa.eu/eli/reg/2024/1689/oj}

\bibitem{opencvzoo}
OPENCV. \textit{OpenCV Zoo: Models for OpenCV DNN} [softvér].
Dostupné z: \url{https://github.com/opencv/opencv_zoo}

\bibitem{minivision2020silentface}
MINIVISION AI. \textit{Silent-Face-Anti-Spoofing} [softvér]. 2020.
Dostupné z: \url{https://github.com/minivision-ai/Silent-Face-Anti-Spoofing}

\bibitem{huang2007lfw}
HUANG, Gary B., RAMESH, Manu, BERG, Tamara a LEARNED-MILLER, Erik.
\textit{Labeled Faces in the Wild: A Database for Studying Face Recognition in
Unconstrained Environments}. Technical Report 07-49. Amherst: University of
Massachusetts, 2007.

\bibitem{zhang2012casia}
ZHANG, Zhiwei, YAN, Junjie, LIU, Sifei, LEI, Zhen, YI, Dong a LI, Stan Z.
A Face Antispoofing Database with Diverse Attacks.
In: \textit{2012 5th IAPR International Conference on Biometrics (ICB)}.
IEEE, 2012, s.~26--31.
DOI: 10.1109/ICB.2012.6199754

\end{thebibliography}

\appendix

\chapter{Obsah repozitára}
\label{app:repo}

Repozitár \url{https://github.com/opoiasnik/face-auth-antispoofing}:

\begin{itemize}
  \item \texttt{backend/} -- serverová aplikácia, testy, vyhodnocovacie skripty,
  \item \texttt{frontend/} -- klientska aplikácia,
  \item \texttt{docs/} -- technická dokumentácia,
  \item \texttt{report/} -- zdrojové súbory tejto správy,
  \item \texttt{docker-compose.yml}, \texttt{render.yaml} -- nasadenie.
\end{itemize}

\chapter{Ukážka odpovede API}
\label{app:api}

\begin{lstlisting}[caption={Zamietnutie relácie pri útoku statickým obrázkom}]
HTTP/1.1 401 Unauthorized
{
  "error": {
    "code": "biometric_rejected",
    "message": "Liveness challenge was not completed",
    "details": {
      "reason": "active_challenge_failed",
      "passive_liveness": 0.9592,
      "active_liveness": {
        "passed": false,
        "steps": [
          {"action": "center",    "passed": true,  "usable_frames": 7,  "peak_delta": 0.1045},
          {"action": "look_up",   "passed": false, "usable_frames": 10, "peak_delta": 0.0},
          {"action": "look_down", "passed": false, "usable_frames": 10, "peak_delta": 0.0}
        ]
      },
      "consistency": 1.0,
      "quality_issues": {},
      "failed_action": "look_up"
    }
  }
}
\end{lstlisting}

\end{document}
